\emph{(i).} With $\lambda$ a field, the linearised constraints are $2\ip{\delta I}{I} - \delta\lambda/3 = 0$, $\ip{\delta V}{I} + \ip{V}{\delta I} = 0$ and $2\ip{\delta V}{V} - 3\,\delta\lambda/\lambda^{2} = 0$, the free case of Proposition \claimref{dark-energy.under-action-written-norm-parameter-laws}: three independent conditions on thirteen components, a ten-dimensional constrained space, inside which the images $(\xi\cdot I, \xi\cdot V, 0)$ of the broken generators span nine dimensions (Proposition \claimref{dark-energy.point-infinity-has-no-physical-fluctuation}: bivectors do not change a norm, and the stabiliser of $(I,V)$ is the six-dimensional Lorentz algebra of $\{I,V\}^{\perp}$) and the norm mode $n$ is constrained and not a gauge direction, so the two together span the constrained space. Gauge transformations act on $(\nabla, I, V)$ and leave $\lambda$ and $T$ alone, so after removing the gauge part the remaining fluctuation of $(I, V, \lambda)$ at each point is $\delta\lambda(x)\,n(x)$. The check recomputes the ranks $3$, $9$, $10$ exactly in the flat model at $\Lambda = 3$, $12$ and $3/4$.

\emph{(ii).} The added term is the only one containing $T$, and $\delta_{T}\,c\int\lambda\,dT = c\int\lambda\,d\,\delta T = \pm c\int d\lambda \wedge \delta T$ up to a boundary term, for every compactly supported $3$-form $\delta T$; so $d\lambda = 0$. On connected $M$, $\lambda = \Lambda$ is constant and $\delta\lambda = \delta\Lambda$ is a constant, so by (i) the only non-gauge fluctuation of $(I, V, \lambda)$ is the global one, $\delta\Lambda\,n$, the tangent at $s = 1$ of the family $(sI, V/s, s^{2}\Lambda)$ up to the factor $1/(2\Lambda)$. Every local fluctuation of $(I, V)$ is therefore gauge, as in Proposition \claimref{dark-energy.point-infinity-has-no-physical-fluctuation}, and its conclusion for the graviton holds unchanged.

\emph{(iii).} On $M = \R \times \Sigma$ write $T = T_{\Sigma} + dt \wedge T_{0}$ with $T_{\Sigma}$ a $3$-form and $T_{0}$ a $2$-form on the slices. Then $dT = dt \wedge (\partial_{t}T_{\Sigma} - d_{\Sigma}T_{0})$ and, integrating by parts on the compact slices, $c\int\lambda\,dT = c\int dt\int_{\Sigma}\bigl(\lambda\,\partial_{t}T_{\Sigma} - d_{\Sigma}\lambda \wedge T_{0}\bigr)$ up to a sign fixed by orientation. So $T_{\Sigma}$ is the momentum conjugate to $\lambda$ (with symplectic form $c\,\delta\lambda \wedge \delta T_{\Sigma}$ integrated over $\Sigma$), and $T_{0}$ is a multiplier for $d_{\Sigma}\lambda = 0$. These constraints are first class (they involve $\lambda$ alone) and the smeared constraint $\int_{\Sigma}d_{\Sigma}\lambda \wedge \theta$ generates $T_{\Sigma} \mapsto T_{\Sigma} \pm d_{\Sigma}\theta$. The reduced space is the constraint surface, $\lambda$ constant on $\Sigma$, that is $H^{0}(\Sigma)$, times the quotient of $3$-forms by exact ones, $H^{3}(\Sigma)$; for $\Sigma$ compact, connected and orientable both are one-dimensional, and the pairing $\int_{\Sigma}\lambda\,T_{\Sigma} = \Lambda\tau$ between them is nondegenerate (Poincar\'e duality). So the sector reduces to one canonical pair. That the remaining gravitational part depends on $\lambda$ does not change this: it contributes to the evolution of $\tau$, not to the constraints. This is the Henneaux-Teitelboim reduction of Fact \claimref{dark-energy.imported-henneaux-teitelboim-promotion}, here applied to the norm. The check builds the cubical cochain complex of the periodic lattice $(\mathbb{Z}/L)^{3}$, $L = 2, 3$, verifies $\delta^{2} = 0$ and the Betti numbers $1, 3, 3, 1$ of the $3$-torus, and that the sum over cubes vanishes on every coboundary, so that the constants and the total are the surviving pair.

\emph{(v).} The reduced symplectic form is $c\,d\Lambda \wedge d\tau$, whose Liouville measure is $|c|\,d\Lambda\,d\tau$: its density per unit $\tau$ is constant in $\Lambda$, and $\int_{0}^{\infty}d\Lambda$ diverges. Turning this into a relative weight over $\Lambda$ needs a convention for the $\tau$-range attached to each value; $\tau$ is non-compact and on a solution advances by a solution-dependent amount, so the equal-range convention used for the fraction below is an assumption, not a consequence. Promoting $g(\lambda)$ instead replaces $\lambda$ by $g(\lambda)$ in (iii), so the symplectic form becomes $c\,g'(\Lambda)\,d\Lambda \wedge d\tau$: density $|c|/3$ for $g = \lambda/3$, $3|c|/\Lambda^{2}$ for $g = 3/\lambda$. With the flat density, equal $\tau$-ranges and support $(0, \Lambda_{\max}]$ the measure of $(0, \Lambda_{\mathrm{obs}}]$ is the fraction $\Lambda_{\mathrm{obs}}/\Lambda_{\max}$, and $N^{2} = 3M_{P}^{2}/\Lambda$ turns it into $3/N^{2}$ for $\Lambda_{\max} = M_{P}^{2}$. The check verifies the two derivatives through exact difference quotients and the last identity at rational points.

\emph{(iv).} Varying $\nabla$, $I$ and $V$ in the promoted functional at constant $\lambda = \Lambda$ gives the equations of the unpromoted functional at parameter $\Lambda$, since the added term contains none of them; varying $\lambda$ gives $c\,dT = -\partial\mathcal{L}/\partial\lambda$, a $4$-form equation on $M = \R \times \Sigma$, which is solvable because a noncompact connected $4$-manifold has $H^{4}(M) = 0$. So each solution of the unpromoted equations at parameter $\Lambda$, which by Proposition \claimref{dark-energy.under-action-written-norm-parameter-laws}(i) has $N = M_{P}\sqrt{3/\Lambda}$, extends to a solution of the promoted functional, and these solution spaces for different $\Lambda$ all lie in the one solution space of the promoted functional. On it $N$ takes different values, so it is not constant there; the promoted theory is a single functional, so by Definition \claimref{obj.dark-energy.four-ways-number-enters-theory-shape} $N$ is (I) and not (P) or (D).
